
O UV mapping associa cada ponto de uma superfície 3D a um ponto de uma textura 2D. Seguindo Heckbert \cite[p.~2]{heckbert1986}, o processo foi tratado como a composição de uma parametrização da superfície com a projeção para a tela. O UV dinâmico altera essa parametrização a cada quadro sem tocar na malha, e o mapeamento procedural substitui a imagem por uma função matemática \cite[p.~1]{dong2020}.

As três abordagens resolvem problemas diferentes. O UV tradicional é a escolha para detalhes únicos e posicionais, como rostos, logotipos e decalques. O UV dinâmico é a escolha para movimento repetitivo (água, lava, esteiras, campos de energia), com o custo de uma amostragem e poucas operações por fragmento. O procedural é a escolha para grandes áreas sem repetição e para famílias de variações \cite[p.~2]{dong2020}. Em troca, ele exige mais cálculo por fragmento, não tem mipmaps que o protejam do aliasing e é difícil de controlar. Dong et al.\ descrevem o ajuste de parâmetros como um processo demorado \cite[p.~23]{dong2020}, e Smelik et al.\ apontam a baixa controlabilidade como o principal motivo de os métodos procedurais ainda não substituírem a modelagem manual \cite[p.~1]{smelik2014}.

A análise das equações mostrou que a pergunta mais importante para o desempenho é onde cada cálculo é executado. Transformações lineares podem ir para a CPU ou para o vertex shader. Deformações não lineares e funções procedurais precisam ficar no fragment shader, que executa centenas de vezes mais que o vertex shader no modelo usado. No Vulkan, a técnica depende de um sampler com endereçamento \texttt{REPEAT}, de um descriptor set que exponha a textura e os parâmetros, e de um mecanismo para enviar o tempo a cada quadro. A implementação usa um Uniform Buffer porque os command buffers são pré-gravados; push constants seriam a alternativa natural se fossem regravados por quadro.

Os produtos analisados mostram as duas formas de uso na indústria. Nos motores de jogo, o UV dinâmico é um recurso de material executado em tempo real, como o \textit{Panner} da Unreal e o \textit{Tiling And Offset} da Unity. Nas ferramentas de autoria, como Substance Designer, SpeedTree e CityEngine, o procedural é usado durante a produção. No Rocket League, as texturas das arenas e das rodas são criadas com o Substance Designer durante a produção, e o jogo as aplica como texturas comuns sobre o UV de cada modelo \cite{substance}. No Marvel's Spider-Man Remastered, a técnica aparece nos trajes, nos decalques e nas fachadas repetidas, e a água da baía é um caso compatível com UV dinâmico. Como os estúdios não publicam a implementação desses efeitos, o UV dinâmico em funcionamento foi demonstrado na implementação do grupo, em que o deslocamento da textura pode ser medido quadro a quadro.

A implementação entregue tem limitações conhecidas, discutidas ao longo do texto: a textura não tem mipmaps, o tempo não é mantido periódico, a translação de UV ainda é feita por pixel e a grade procedural não tem antialiasing analítico. Esses pontos são o caminho natural para uma próxima versão, junto com a divisão dos parâmetros entre Uniform Buffer e push constants.
